\documentclass[9pt,letterpaper]{article}
\usepackage[margin=0.60in]{geometry}
\usepackage{graphicx,booktabs,tabularx,array,xcolor,amsmath,hyperref}
\usepackage{enumitem}
\usepackage{caption}
\usepackage{microtype}
\hypersetup{colorlinks=true,urlcolor=blue,linkcolor=blue}
\setlength{\parindent}{0pt}
\setlength{\parskip}{2pt}
\setlist[itemize]{nosep,leftmargin=*}
\captionsetup{font=small,labelfont=bf}
\begin{document}
\begin{center}
{\Large\bfseries Distance-3 Surface-Code Memory on IQM Emerald}\\[4pt]
{\large Internal Experimental Summary}\\[3pt]
{\small DuQuantum IQM Challenge 2026 \;|\; 24 October 2026}
\end{center}
\vspace{-6pt}
\section*{Objective and main finding}
We evaluated whether more frequent stabilizer extraction improves the storage of logical $|0_L\rangle$ (Z-memory) and $|+_L\rangle$ (X-memory) on IQM Emerald, while holding the \emph{memory interval} approximately fixed at $15\,\mu\mathrm{s}$. The measured implementation showed \textbf{no net benefit from additional intermediate rounds}: Z-memory became worse as the number of rounds increased, while X-memory remained close to a 50\% logical error probability. Unencoded single-qubit memory performed better under the tested conditions. 

\section*{Implementation and experimental design}
The experiment uses a rotated $[[9,1,3]]$ surface code with nine data qubits and eight ancillas on a fixed 17-qubit hardware patch. A common, unitary \emph{known-state} preparation was used for each logical basis; it is not fault-tolerant preparation or an encoder for an arbitrary unknown input state. The qubit mapping was optimized and verified to require no additional routing SWAPs at the native-gate level. Intermediate ancilla resets were omitted. Measurement histories were converted to detection events, then decoded with PyMatching minimum-weight perfect matching (MWPM) using an arm-specific Stim detector error model. No physical correction gates were applied during storage. We also tried tuning decoder parameters, but the improvement was limited; alternative decoders remain an interesting direction to explore.

The $15\,\mu\mathrm{s}$ interval starts after preparation and ends before final basis rotation / destructive readout; it includes all intermediate extraction operations and idle delays. For $r=0,1,2,3$ rounds, available idle time was distributed approximately uniformly over $r+1$ intervals. Locally compiled memory windows were $15.000\,\mu\mathrm{s}$ except for $r=1,2$ ($14.996\,\mu\mathrm{s}$). Timings were extracted from IQM pulse-compiler schedules rather than summed nominal gate lengths. Single-qubit $|0\rangle$ and $|+\rangle$ controls used the same defined memory interval. Each of 14 configurations was measured in two blocks of 1,000 shots, for \textbf{28,000 total shots}. All arms shared the saved calibration ID.

\section*{Measured results}
\begin{table}[h]
\centering\small
\begin{tabular}{lrr}
\toprule
\textbf{15 $\boldsymbol{\mu\mathrm{s}}$ memory configuration} & \textbf{Z error (95\% CI)} & \textbf{X error (95\% CI)}\\
\midrule
Encoded, 0 intermediate rounds & 25.10\% (23.25--27.05) & 48.50\% (46.31--50.69)\\
Encoded, 1 intermediate round & 32.30\% (30.29--34.38) & 45.45\% (43.28--47.64)\\
Encoded, 2 intermediate rounds & 36.50\% (34.42--38.63) & 49.90\% (47.71--52.09)\\
Encoded, 3 intermediate rounds & 37.40\% (35.31--39.54) & 48.60\% (46.41--50.79)\\
Unencoded single qubit & \textbf{0.65\%} (0.38--1.11) & \textbf{29.55\%} (27.59--31.59)\\
\bottomrule
\end{tabular}
\caption{Logical failure rates after decoding (encoded arms) and measured state error (unencoded controls); 2,000 shots per arm. Intervals are Wilson 95\% binomial intervals, shown in percentage points.}
\end{table}
\vspace{-10pt}
\begin{center}
\includegraphics[width=0.56\linewidth]{memory_15us.png}
\end{center}
\vspace{-15pt}
\section*{Zero-round storage diagnostic}
To distinguish end-to-end preparation/readout error from storage-related degradation, we also ran zero-intermediate-round encoded memories at near-zero and $5\,\mu\mathrm{s}$ storage. The same complete state-preparation and final-readout procedures were retained.
\begin{table}[h]
\centering\small
\begin{tabular}{lrr}
\toprule
\textbf{Encoded, zero intermediate rounds} & \textbf{Z error} & \textbf{X error}\\
\midrule
Near-zero memory interval & 2.90\% & 2.00\%\\
$5\,\mu\mathrm{s}$ & 11.70\% & 52.25\%\\
$15\,\mu\mathrm{s}$ & 25.10\% & 48.50\%\\
\bottomrule
\end{tabular}
\caption{Each entry represents 2,000 hardware shots. X-memory at 5 and 15 $\mu$s is close to random logical-parity prediction.}
\end{table}
\vspace{-10pt}
\begin{center}
\includegraphics[width=0.56\linewidth]{zero_round_storage.png}
\end{center}
\vspace{-15pt}
\section*{Component-control experiment (separate Z-memory run)}
We also compared a one-round baseline against added CZ echoes, repeated ancilla measurements, and a full three-round extraction circuit. These controls came from an earlier Z-memory experiment with a different preparation protocol; they should be compared \emph{within this experiment}, not pooled with the fixed-$15\,\mu$s benchmark above. The locally compiled full schedules were matched within $4\,$ns, including preparation and readout.
\begin{table}[h]
\centering\small
\begin{tabular}{lrrr}
\toprule
\textbf{Circuit} & \textbf{Decoded LER (95\% CI)} & \textbf{CZ count} & \textbf{Shots}\\
\midrule
1 round + idle & 18.75\% (17.10--20.52) & 19 & 2,000\\
1 round + CZ echo & 19.70\% (18.02--21.50) & 67 & 2,000\\
1 round + ancilla measurements & 18.30\% (16.67--20.05) & 19 & 2,000\\
3 complete rounds & \textbf{32.55\%} (30.53--34.64) & 67 & 2,000\\
\bottomrule
\end{tabular}
\caption{Hardware component controls; 2,000 shots per arm. Extra CZ gates were paired echoes; auxiliary ancilla measurements were not included in the decoder.}
\end{table}
The three-round LER exceeded the one-round-plus-idle baseline by \textbf{13.80 percentage points} (unadjusted 95\% CI: 11.12--16.46). Neither isolated CZ echoes (+0.95 pp) nor auxiliary ancilla measurements ($-0.45$ pp) showed a detectable difference from the idle control. Thus, the full extraction sequence is substantially more damaging than either tested intervention alone; these controls do not uniquely identify the underlying noise mechanism.

\section*{Discussion and takeaway}
At fixed $15\,\mu$s storage, added extraction rounds increased Z-memory LER but did not improve X-memory. In the separate schedule-matched component controls, full three-round extraction raised Z error by 13.80 percentage points relative to one round plus idle, whereas the tested isolated CZ-echo and ancilla-measurement additions produced no detectable changes. The component tests do not isolate a unique physical noise source. Near-zero storage errors of only 2--3\% point to substantial memory-interval degradation. Overall, the tested encoded memories did not beat unencoded qubits. Zero-SWAP mapping was verified and decoder parameter tuning gave limited improvements; other decoders and noise models remain worth exploring.

\vfill
{\footnotesize\color{gray}Source: provided \texttt{submission(1).zip}, especially \texttt{METHODS.md}, \texttt{RESULTS.md}, \texttt{validation.json}, and archived measurement tables. Figures are reproduced from the submission package.}
\end{document}
